The term \QSOpen{}vacuum\QSClose{} has distinct meanings across classical physics, quantum field theory, and general relativity. In classical physics, vacuum refers to the absence of material particles—a region of space devoid of atoms, molecules, or macroscopic matter. The classical vacuum is an empty stage on which forces act.

In quantum field theory (QFT), the notion of vacuum acquires a different character. Here, fields—not particles—constitute the primary entities. The vacuum is the ground state of all quantum fields, the configuration of lowest possible energy consistent with the commutation relations and field dynamics. Even when no particles are present, quantum fields fluctuate around their minima, giving rise to nonzero vacuum expectation values for certain observables. These fluctuations are not artefacts of measurement or disturbance; they are features of the quantum equations themselves. In Minkowski spacetime, the vacuum is Lorentz invariant. No preferred direction or frame exists, and the absence of particles is an absolute property shared by all inertial observers.

However, in general relativity (GR), spacetime is no longer a fixed, flat background. It becomes a dynamical entity whose curvature interacts with matter and energy. The introduction of curved spacetime disrupts the global symmetries that underlie the inertial vacuum of QFT. In regions of strong gravitational fields or global curvature, there is generally no unique, globally defined vacuum state. Instead, the concept of vacuum becomes observer-dependent. Different families of observers may disagree about whether a given region of spacetime is populated by particles. This relativity of the vacuum holds because the definition of positive frequency modes—those corresponding to particle excitations—depends on the choice of time coordinate, which itself is tied to the observer's worldline. Consequently, what one observer identifies as an empty vacuum, another observer may interpret as a state containing particles, momentum, or thermal radiation.

In quantum field theory, the notion of a particle is defined relative to a choice of time. In flat Minkowski spacetime, all inertial observers share the same time direction (up to Lorentz boosts that preserve the sign of frequency). This gives a unique way to split field solutions into positive-frequency modes (time dependence $e^{-i\omega t}$, $\omega > 0$) and negative-frequency modes. Creation and annihilation operators are built from this split, and the vacuum is the state annihilated by all annihilation operators. Because inertial observers agree on which frequencies are positive, they agree on particle number.

In curved spacetimes or for accelerating observers, no universal time direction exists. Different observers, following different trajectories, adopt different time coordinates and decompose the field into different sets of modes. What one observer identifies as a positive-frequency mode may contain negative-frequency components according to another. Since positive frequency defines particles and negative frequency defines antiparticles, the two observers disagree on how many particles are present.

Mathematically, if one observer expands the field $\hat{\phi}(x)$ in terms of a basis of modes $\{f_i(x)\}$, while another observer uses a different basis $\{g_j(x)\}$, the two expansions are related by a Bogoliubov transformation. This transformation mixes creation and annihilation operators, leading to the possibility that the vacuum state for one observer appears populated with particles to another. If the Bogoliubov coefficients $\beta_{ij}$ are nonzero, the expected particle number for mode $j$ in the $g$-basis, measured in the $f$-vacuum, is $\langle 0_f | \hat{N}_{g,j} | 0_f \rangle = \sum_i |\beta_{ij}|^2$, and the total count is $\langle 0_f | \hat{N}_g | 0_f \rangle = \sum_{i,j} |\beta_{ij}|^2$.

The number operator $\hat{N} = \sum_j \hat{a}_j^\dagger \hat{a}_j$ counts excitations above a given vacuum. It is constructed from creation and annihilation operators, which themselves depend on the choice of mode basis. Change the basis—by switching to an accelerating frame or a different time slicing—and the operator changes. Two observers examining the same quantum state with their respective number operators obtain different expectation values. This is not a coordinate artefact—a particle detector carried by an accelerating observer registers real energy deposits while an inertial detector in the same region remains silent.

A particle, in this framework, is not a fundamental object embedded in spacetime. The quantum field is fundamental—it pervades all of spacetime and its state is agreed upon by all observers. A particle is what a detector registers when coupled to the field along a specific worldline. The Unruh-DeWitt model makes this precise. It describes a pointlike system with discrete energy levels that follows a trajectory through spacetime, interacting locally with the field at each point along its path. The excitation rate depends entirely on the trajectory. Along an inertial path in the Minkowski vacuum, the detector remains in its ground state, but along a uniformly accelerating path through the same field state, it absorbs energy and clicks. The same field, the same quantum state, two different particle counts.

Classical wave physics offers a partial analogy. A moving listener hears a Doppler-shifted frequency, but the wave carries the same energy—no new physics follows from the listener's motion. In quantum field theory, frequency decomposition is more consequential. Positive-frequency modes define particles; negative-frequency modes define antiparticles. When an observer's trajectory mixes these—as acceleration does—modes that one observer calls vacuum become, for the other, a superposition containing particles. The accelerating detector absorbs real energy, thermalises, and reaches equilibrium at a definite temperature. This energy is drawn from the agency maintaining the acceleration, so the total energy budget remains consistent across frames. The field state has not changed; what has changed is which degrees of freedom the detector accesses.

Yet the quantum state of the field is observer-independent. The state vector $|\Psi\rangle$ in the Hilbert space does not change when observers switch coordinates. What changes is the decomposition of that state into particle and vacuum sectors. The stress-energy tensor $T_{\mu\nu}$, local curvature invariants, and S-matrix elements between asymptotic states also remain invariant. The field equations hold in all frames. What becomes observer-dependent is the particle interpretation—how many quanta a given detector counts—not the underlying field configuration or its gravitational effects.

The Unruh effect (Unruh, 1976) occurs in flat Minkowski spacetime. An inertial observer in this vacuum detects nothing, but an observer undergoing uniform acceleration $a$ through the same region perceives a thermal bath at temperature $T_U = \hbar a/(2\pi c k_B)$. The accelerating observer follows hyperbolic trajectories that asymptotically approach but never cross a Rindler horizon—a boundary in flat spacetime beyond which signals can never reach the accelerating observer, analogous to a black hole horizon. Rindler coordinates, natural to this observer, cover only the wedge of spacetime accessible to them. The Minkowski vacuum, expressed in Rindler modes, decomposes into entangled pairs straddling this horizon. Since the accelerating observer cannot access modes behind the horizon, tracing over those inaccessible modes produces a thermal density matrix. The thermal spectrum is exact for uniform acceleration, not an approximation.

Hawking radiation (Hawking, 1974) arises from an analogous mechanism at black hole horizons. An observer falling freely across the horizon encounters no particles—locally, the equivalence principle guarantees that spacetime looks flat, and the vacuum looks empty. A distant stationary observer, however, perceives thermal radiation at temperature $T_H = \hbar c^3/(8\pi G M k_B)$, where $M$ is the black hole mass. The event horizon separates spacetime into causally disconnected regions. Field modes straddling the horizon are entangled across it. The distant observer, who can access only the exterior, must trace over the interior modes. The resulting reduced state is thermal, with its spectrum set by the surface gravity. The connection to the Unruh effect is direct: a stationary observer hovering near a black hole must accelerate to avoid falling in, and the temperature they measure matches the Unruh temperature for that acceleration. At large distance, the gravitational redshift reduces this to the Hawking temperature.

Cosmological particle production occurs in expanding spacetimes. During inflation, the rapid exponential expansion of space stretches quantum vacuum fluctuations to macroscopic scales. A mode that begins as a short-wavelength fluctuation well inside the Hubble radius gets stretched until its wavelength exceeds the horizon size. At that point the mode freezes—it can no longer oscillate because no causal process spans its wavelength—and its quantum amplitude becomes a classical perturbation in the energy density. After inflation ends and the expansion decelerates, these frozen modes re-enter the horizon and seed the density perturbations that grow into galaxies and galaxy clusters. The temperature anisotropies in the cosmic microwave background, measured by COBE (1992), WMAP (2003), and Planck (2013), match the nearly scale-invariant spectrum predicted by this mechanism. The particle interpretation follows the same pattern as the Unruh and Hawking cases: the rapid change in the metric mixes positive and negative frequencies, and an observer at a later epoch counts particles where an earlier observer counted none.

The Casimir effect (Casimir, 1948) demonstrates that boundary conditions alter the vacuum. Two parallel conducting plates restrict which field modes can exist between them, lowering the vacuum energy density relative to the space outside. The resulting pressure difference pushes the plates together—a measurable force produced by the vacuum itself. In the dynamical variant, the plates accelerate or their separation oscillates, and the changing boundary conditions parametrically amplify vacuum fluctuations into real photons. Laboratory realisations using superconducting circuits with time-modulated boundaries have confirmed this photon creation.

\begin{commentary}[Hawking Radiation]
By 1972, black hole mechanics followed formal laws resembling thermodynamics: changes in mass, angular momentum, and charge obeyed a relation involving surface gravity $\kappa$ and horizon area $A$, with $\kappa$ positioned where temperature would appear and $A$ where entropy belongs. Yet classical general relativity permitted no emission, implying zero temperature. The analogy appeared to be coincidence.

Jacob Bekenstein argued otherwise. On information-theoretic grounds, black holes must possess entropy proportional to horizon area. Hawking rejected this—without an emission mechanism, a black hole could not have true temperature.

In 1973, Hawking visited Moscow and met Yakov Zel'dovich and Alexei Starobinsky, who had shown that rotating bodies could amplify vacuum fluctuations—superradiance. Hawking set out to prove gravitational collapse could not produce particles.

He traced field modes through the spacetime of a collapsing star. Outgoing modes near the forming horizon undergo exponential redshift, mixing positive and negative frequencies. The Bogoliubov coefficients yield a Planck spectrum at temperature $T_H = \hbar c^3/(8\pi G M k_B)$. Hawking spent weeks rechecking, yet the result persisted.

Hawking published in \textit{Nature} as \DSQOpen{}Black hole explosions?\DSQClose{} (1974), with a full derivation in \textit{Communications in Mathematical Physics} (1975). The result introduced the information paradox: thermal radiation carries no trace of what fell in, contradicting unitarity.

The popular picture of virtual particle pairs splitting at the horizon does not appear in Hawking's derivation. The radiation follows from the global causal structure of the collapsing spacetime and the entanglement of field modes across the horizon. It has never been directly observed—for a solar-mass black hole, $T_H \approx 6 \times 10^{-8}$ K—though analogue experiments in BECs, fluids, and superconducting circuits have reproduced consistent features.

\end{commentary}

\vspace*{\fill}
\begin{table}[ht]
\centering
\footnotesize
\renewcommand{\arraystretch}{1.3}
\setlength{\tabcolsep}{3pt}
\begin{tabular}{|>{\centering\arraybackslash}m{2.0cm}|
                >{\centering\arraybackslash}m{2.2cm}|
                >{\centering\arraybackslash}m{2.2cm}|
                >{\centering\arraybackslash}m{2.2cm}|
                >{\centering\arraybackslash}m{2.2cm}|
                >{\centering\arraybackslash}m{2.2cm}|}
\hline
\textbf{Feature} & \textbf{Unruh Effect} & \textbf{Hawking Radiation} & \textbf{Cosmological Particle Creation} & \textbf{Schwinger Effect} & \textbf{Dynamical Casimir Effect} \\
\hline
\textbf{Physical Context} & Uniform acceleration in flat spacetime & Black hole event horizon & Expanding FLRW universe & Strong external electric field & Time-dependent boundary conditions \\
\hline
\textbf{Primary Mechanism} & Bogoliubov transformation (Minkowski $\leftrightarrow$ Rindler) & Horizon mode mismatch and equivalence principle & Mode stretching and horizon crossing & Vacuum instability via tunnelling & Parametric amplification of vacuum fluctuations \\
\hline
\textbf{Key Dependence} & Proper acceleration $a$ & Black hole mass $M$, surface gravity & Hubble rate $H$, coupling strength & Electric field strength $E$ & Modulation frequency and boundary speed \\
\hline
\textbf{Characteristic Scale / Formula} & $T_U = \dfrac{\hbar a}{2\pi c k_B}$ & $T_H = \dfrac{\hbar c^3}{8\pi G M k_B}$ & Particle density $\propto H^2$ & $\Gamma \propto \exp\left(-\dfrac{\pi E_c}{E}\right)$ & Photon production peaks at $\omega_{\text{mod}} \approx 2\omega_{\text{cav}}$ \\
\hline
\textbf{Experimental Probes} & Analogues: BECs, trapped ions, superconducting circuits & Analogues (BECs, fluids); black hole thermodynamics & CMB anisotropies, primordial fluctuations & High-intensity lasers, graphene lattices & SQUID-based circuits, modulated cavities \\
\hline
\end{tabular}
\caption*{
\centering
Comparative overview of observer-dependent vacuum phenomena. Each column represents a distinct physical context in which the concept of vacuum, and thus of particle content, becomes relative to the observer's state of motion or horizon access.
}
\end{table}
\vspace*{\fill}



